% =====================================================================
%  ThmSmith Stage -1 proposal skeleton.
%  Written automatically by the ThmSmith TS pipeline before Stage -1 runs.
%  Locked parameters: qid=eid\_shift\_share\_design, specialization=v1
%
%  HOW TO USE THIS FILE
%  --------------------
%  - The preamble (everything above the `% --- BEGIN CONTENT ---` marker)
%    and the `\section{...}` headers below are fixed by the pipeline.
%    Do NOT rewrite or rename them; the Step 3 reviewer subagent and the
%    downstream Stage 0 / Stage 1 prompts grep for these section titles.
%  - Each section contains exactly one `% TODO(stage-1 ...): ...`
%    placeholder. Replace each TODO with the content described for that
%    section in `ThmSmith/tools/src/prompts/stage_neg1_draft.txt`
%    ("REQUIRED OUTPUT FORMAT (Stage -1 proposal .tex)").
%  - The `\paragraph{Locked parameters.}` block in the metadata block
%    must pin the chosen `(qid, specialization, p, assignment, target,
%    regression)` precisely. If the proposal maps to the Q1 pool-mode,
%    reproduce that block verbatim from
%    `ThmSmith/doc/panel_formula_question_generator_note_with_common_prompts.tex`
%    (Q2..Q6 templates have been retired). Otherwise (free-form
%    `(qid, specialization)` — the default), define each parameter
%    explicitly here (no implicit conventions). Stage 0 quotes this block;
%    do not rephrase it after locking.
%  - Removing a section header, renaming a macro, or rewriting the
%    preamble is a regression: the file must still match this skeleton
%    after you finish.
% =====================================================================

\documentclass[11pt]{article}

\usepackage{amsmath,amssymb,amsthm,mathtools}
\usepackage{enumitem}
\usepackage[colorlinks=true,linkcolor=blue,citecolor=blue,urlcolor=blue]{hyperref}
\usepackage{natbib}

\newcommand{\E}{\mathbb{E}}
\renewcommand{\Pr}{\mathbb{P}}
\newcommand{\one}{\mathbf{1}}
\newcommand{\transpose}{\mathsf{T}}
\newcommand{\calH}{\mathcal{H}}
\newcommand{\calF}{\mathcal{F}}
\newcommand{\calR}{\mathcal{R}}
\newcommand{\R}{\mathbb{R}}
\newcommand{\plim}{\operatorname*{plim}}

\theoremstyle{definition}
\newtheorem{definition}{Definition}
\newtheorem{assumption}{Assumption}
\newtheorem{example}{Example}

\theoremstyle{plain}
\newtheorem{theorem}{Theorem}
\newtheorem{conjecture}{Conjecture}
\newtheorem{proposition}{Proposition}
\newtheorem{lemma}{Lemma}
\newtheorem{corollary}{Corollary}

\theoremstyle{remark}
\newtheorem{remark}{Remark}

\title{ThmSmith Stage -1 proposal: \texttt{eid\_shift\_share\_design} / \texttt{v1}%
  \\ \large\textit{Leakage-null projection frontier for recentered and design-based formula IV}}
\author{ThmSmith Stage -1 proposer}
\date{\today}

\begin{document}
\maketitle

% --- BEGIN CONTENT ---  (everything above is fixed by the pipeline)

\paragraph{Locked parameters.}
\begin{verbatim}
qid = eid_shift_share_design
specialization = v1
anchor cluster = Exact identification / formula-IV shift-share designs
graph / SWIG = a linear structural IV SWIG with predetermined exposure
  shares S, observed formula inputs W, shock vector G, treatment X(S,G),
  outcome Y, structural exposure coefficient beta_exp, and outcome error U.
treatment = scalar exposure-driven treatment X_i generated by a known
  formula from shares S_i and shocks G.
target estimand = beta_exp in E[Y_i | X_i,U_i] = beta_exp X_i + U_i,
  equivalently the constant structural exposure coefficient identified by
  valid formula-IV moments.
identifying assumptions = linear exposure additivity, consistency, positive
  first stage for every retained shock filter, predetermined shares, known
  counterfactual shock design, and the projection-leakage restrictions stated
  in Assumptions 1--7 below.
canonical estimator = Bartik/SSIV moment using the unfiltered shock vector G.
recentered estimator = Borusyak-Hull recentered moment using the filter
  A_R = I - P_R, where P_R projects onto the span removed by counterfactual
  shock-mean recentering.
design estimator = design-based formula-IV moment using the filter
  A_D = I - P_D, where P_D projects onto the larger covariate-conditional
  design span, with Range(P_R) subset Range(P_D).
novel kernel = characterize when epsilon is killed by exposure aggregation
  after A_R or A_D filtering: canonical, recentered, and design-IV consistency
  are governed by Pi^{1/2} epsilon, Pi^{1/2} A_R epsilon, and
  Pi^{1/2} A_D epsilon, respectively.
\end{verbatim}

\begin{abstract}
Formula-IV papers now give three adjacent routes for shift-share identification: canonical shock-level SSIV, recentered instruments for nonrandom exposure, and design-based formula IV that conditions on a known shock assignment process.  This proposal asks when the latter two are genuinely different rather than notational variants on the same population moment.  The proposed kernel is a leakage-null projection theorem: after exposure aggregation, canonical Bartik, recentered IV, and design IV are consistent exactly when the leakage vector \(\epsilon\) lies in three nested null spaces, \(\ker(\Pi^{1/2})\), \(\ker(\Pi^{1/2}A_R)\), and \(\ker(\Pi^{1/2}A_D)\).  The flagship conjecture is that the middle and largest null spaces are sharply separated on an open class of finite-support formula designs, so design conditioning can identify \(\beta_{\exp}\) in cases where recentering still converges to a biased limit.  Resolving this would give practitioners a finite-support matrix diagnostic for whether recentering has actually removed the harmful covariate-conditional shock-mean leakage.
\end{abstract}

\section{Introduction}
\paragraph{Why the problem is important.}
Shift-share and other formula instruments are often defended by saying that the shocks are as-good-as-random after the right adjustment.  In the current literature, however, the adjustment differs across canonical SSIV, recentered instruments, and design-based formula IV.  The methodological question is whether these adjustments are alternative implementations of the same identifying moment, or whether there is a sharp class of nonrandom-exposure designs in which they identify different population functionals.

\paragraph{The main finding (proposed).}
The proposed headline is Conjecture~\ref{conj:leakage-null-frontier}: in a finite-support linear formula-IV design, the three estimator families have nested leakage-null regions.  Canonical Bartik requires \(\Pi^{1/2}\epsilon=0\), Borusyak-Hull recentering requires \(\Pi^{1/2}A_R\epsilon=0\), and design-based formula IV requires \(\Pi^{1/2}A_D\epsilon=0\).  Since \(A_D\) removes a larger covariate-conditional design span than \(A_R\), there is an open strict-extension region where design IV identifies \(\beta_{\exp}\) but recentered IV and canonical SSIV do not.

\paragraph{What is new.}
The novelty kernel is not another sufficient validity condition for shift-share IV.  The closest paper, \citet{BorusyakHull2023Econometrica}, shows how to recenter a formula instrument by subtracting the counterfactual average induced by nonrandom exposure, while \citet{BorusyakHullJaravel2025ECTJ} reviews design-based formula-IV moments.  The epsilon-gap is the proposed iff projection diagnostic: the harmful leakage directions are classified by the observable exposure covariance \(\Pi\) and by the two removal operators \(A_R,A_D\), rather than by a single statement that a design expectation has been subtracted.

\paragraph{How it positions in the literature.}
This is an exact-identification proposal in the formula-IV and shift-share cluster.  \citet{BorusyakHullJaravel2022REStud} provide the canonical shock-level SSIV equivalence, \citet{BorusyakHull2023Econometrica} provide the recentered-instrument route for nonrandom exposure, and \citet{BorusyakHullJaravel2025ECTJ} organize the design-based formula-IV frontier.  The present proposal asks for the missing equality-versus-strict-extension theorem between these routes, using \citet{GoldsmithPinkhamSorkinSwift2020AER} and \citet{AdaoKolesarMorales2019QJE} as reminders that share geometry can change either the interpretation or the inference even when the instrument has the same visible formula.

\paragraph{Implications for the community.}
The concrete downstream consumer is the diagnostic program in \citet{BorusyakHullJaravel2025ECTJ}: a finite-support \(\Pi,A_R,A_D,\epsilon\) matrix check would tell applied researchers whether the Borusyak-Hull recentered instrument is enough, or whether the richer design-conditioning moment is needed before estimating market-access, trade-shock, or immigration-style formula-IV effects.

\paragraph{How research benefits from the conclusion.}
Readers should report the two projection residual norms \(\|\Pi^{1/2}A_R\epsilon\|\) and \(\|\Pi^{1/2}A_D\epsilon\|\) before treating recentered and design-based formula IV as interchangeable.

\section{Related work}
\citet{BorusyakHullJaravel2022REStud} are the canonical SSIV anchor: they show that shift-share orthogonality can be represented at the shock level and motivate shock-level consistency conditions while allowing endogenous exposure shares.  \citet{BorusyakHull2023Econometrica} are the closest published comparator for this proposal because their recentered-instrument construction subtracts the counterfactual average treatment or instrument generated by nonrandom exposure to shocks.  \citet{BorusyakHullJaravel2025ECTJ} review design-based formula instruments, including known shock designs, recentering, heterogeneous effects, and efficient construction; the proposed projection-null theorem is positioned as a missing sharp comparison inside that review's frontier.  \citet{GoldsmithPinkhamSorkinSwift2020AER} give the exposure-share and Rotemberg-weight interpretation of Bartik instruments, which is relevant because \(\Pi\)-loaded leakage can be invisible under a share-exogeneity narrative but active under shock-level aggregation.  \citet{AdaoKolesarMorales2019QJE} show that shift-share exposure geometry creates serious correlation and inference issues; their contribution is inference rather than the identification-null-space classification proposed here.  \citet{JaegerRuistStuhler2018NBER} provide a substantive shift-share warning in immigration designs, where persistent settlement shares can make a conventional instrument capture long-run adjustment rather than the intended short-run effect.

The in-catalogue duplicate check uses the Stage -1.1 prior-proposal map and \texttt{doc/API.md}.  \texttt{doc/API.md} reports that \texttt{ThmSmith/ExactID} and \texttt{ThmSmith/PartialID} were empty at seed time, so no catalogue entry already states a shift-share/formula-IV theorem.  The nearest banked proposal is \texttt{eid\_lp\_svar\_nonequiv\_v1}, whose headline concerns LP/SVAR row-space non-equivalence rather than formula-IV leakage.  \texttt{q1\_spectral\_phase\_transition\_p1\_markov} is relevant only as a caution that generic spectral thresholds can fail without a primitive observable diagnostic.  \texttt{eid\_dml\_orth\_uniform\_v1} is the closest finite-support covariance/projection diagnostic, but it concerns DML/LRO/ADML variance maps rather than recentered and design-based shift-share IV.

\section{Formal setup}
There are \(N\) observations and \(K\) shock cells.  Let \(S\in\R^{N\times K}\) be the predetermined exposure-share matrix with row \(s_i^\transpose\), let \(W\) denote observed formula inputs and controls, and let \(G\in\R^K\) be the realized shock vector drawn from a known counterfactual design conditional on \(W\).  The scalar treatment is \(X_i=x_i(S_i,G,W_i)\), and the outcome equation is
\[
  Y_i=\beta_{\exp}X_i+U_i,
\]
where \(\beta_{\exp}\) is the target structural exposure coefficient.  The shock-level aggregation of the structural error is \(r:=N^{-1}S^\transpose U\in\R^K\), and the exposure-share covariance entering the formula-IV moment is \(\Pi:=N^{-1}S^\transpose M_W S\), where \(M_W\) residualizes the observation-level controls used in all three estimators.

Let \(\Omega:=\operatorname{Var}(G\mid W)\) be the shock-design covariance.  The leakage vector \(\epsilon\in\R^K\) is defined by the linear reduced form
\[
  \E[r\mid W,G]-\E[r\mid W]=\Omega^{1/2}\epsilon,
\]
after choosing the Moore--Penrose square root on the support of \(\Omega\).  Equivalently, \(\epsilon\) is the covariate-conditional shock-mean error direction that would be zero under ideal shock exogeneity.  Let \(\mathcal R\subseteq\R^K\) be the recentering span removed by the Borusyak-Hull counterfactual-shock mean, and let \(\mathcal D\subseteq\R^K\) be the larger design-conditioning span removed by the formula-IV design moment, with \(\mathcal R\subseteq\mathcal D\).  Denote the orthogonal projections by \(P_R,P_D\), and set \(A_R=I_K-P_R\) and \(A_D=I_K-P_D\).  The canonical, recentered, and design-filtered shock instruments use \(A_B=I_K\), \(A_R\), and \(A_D\), respectively.

The population IV limit for a shock filter \(A\in\{A_B,A_R,A_D\}\) is written
\[
  \beta(A)=\beta_{\exp}+\frac{\Delta(A,\epsilon)}{\Gamma(A)},
  \qquad
  \Delta(A,\epsilon):=\langle \Pi^{1/2}A\epsilon,q_A\rangle,
\]
where \(q_A\in\R^K\) is the filtered first-stage score and \(\Gamma(A)>0\) is the corresponding first-stage denominator.  The diagnostic residuals are
\[
  L_B(\epsilon)=\Pi^{1/2}\epsilon,\qquad
  L_R(\epsilon)=\Pi^{1/2}A_R\epsilon,\qquad
  L_D(\epsilon)=\Pi^{1/2}A_D\epsilon.
\]
The proposed sharp frontier is about the null spaces of these three maps.

\begin{verbatim}
qid = eid_shift_share_design
specialization = v1
anchor cluster = Exact identification / formula-IV shift-share designs
graph / SWIG = a linear structural IV SWIG with predetermined exposure
  shares S, observed formula inputs W, shock vector G, treatment X(S,G),
  outcome Y, structural exposure coefficient beta_exp, and outcome error U.
treatment = scalar exposure-driven treatment X_i generated by a known
  formula from shares S_i and shocks G.
target estimand = beta_exp in E[Y_i | X_i,U_i] = beta_exp X_i + U_i,
  equivalently the constant structural exposure coefficient identified by
  valid formula-IV moments.
identifying assumptions = linear exposure additivity, consistency, positive
  first stage for every retained shock filter, predetermined shares, known
  counterfactual shock design, and the projection-leakage restrictions stated
  in Assumptions 1--7 below.
canonical estimator = Bartik/SSIV moment using the unfiltered shock vector G.
recentered estimator = Borusyak-Hull recentered moment using the filter
  A_R = I - P_R, where P_R projects onto the span removed by counterfactual
  shock-mean recentering.
design estimator = design-based formula-IV moment using the filter
  A_D = I - P_D, where P_D projects onto the larger covariate-conditional
  design span, with Range(P_R) subset Range(P_D).
novel kernel = characterize when epsilon is killed by exposure aggregation
  after A_R or A_D filtering: canonical, recentered, and design-IV consistency
  are governed by Pi^{1/2} epsilon, Pi^{1/2} A_R epsilon, and
  Pi^{1/2} A_D epsilon, respectively.
\end{verbatim}

\section{Assumptions}
\begin{assumption}[Finite linear formula design]\label{ass:finite-linear}
The design has finite \(K\), predetermined \(S\), observed formula inputs \(W\), and a known conditional shock design for \(G\).  The structural equation is \(Y_i=\beta_{\exp}X_i+U_i\), and \(X_i\) is generated by a known additive formula from \(S_i,G,W_i\).
\end{assumption}

\begin{assumption}[Common residualized exposure covariance]\label{ass:pi}
All three estimators use the same observation-level control residualizer \(M_W\), and \(\Pi=N^{-1}S^\transpose M_W S\) is symmetric positive semidefinite with a fixed Moore--Penrose square root \(\Pi^{1/2}\).
\end{assumption}

\begin{assumption}[Known recentering and design spans]\label{ass:spans}
The recentering span \(\mathcal R\) and the design-conditioning span \(\mathcal D\) are linear subspaces of \(\R^K\), \(\mathcal R\subseteq\mathcal D\), and \(P_R,P_D\) are their orthogonal projections.  Define \(A_B=I_K\), \(A_R=I_K-P_R\), and \(A_D=I_K-P_D\).  The filters are exposure-compatible: \(\ker(\Pi^{1/2})\subseteq\ker(\Pi^{1/2}A_R)\subseteq\ker(\Pi^{1/2}A_D)\).
\end{assumption}

\begin{assumption}[Linear leakage representation]\label{ass:leakage}
There exists \(\epsilon\in\R^K\) such that the conditional shock-error leakage in the shock-level residual \(r=N^{-1}S^\transpose U\) is represented by \(\E[r\mid W,G]-\E[r\mid W]=\Omega^{1/2}\epsilon\) on the support of \(\Omega=\operatorname{Var}(G\mid W)\).
\end{assumption}

\begin{assumption}[Filter-specific moment expansion]\label{ass:moment-expansion}
For every \(A\in\{A_B,A_R,A_D\}\), the population IV limit exists and satisfies \(\beta(A)=\beta_{\exp}+\Delta(A,\epsilon)/\Gamma(A)\), with \(\Delta(A,\epsilon)=\langle \Pi^{1/2}A\epsilon,q_A\rangle\) for a first-stage score \(q_A\in\R^K\).
\end{assumption}

\begin{assumption}[Positive first stages]\label{ass:first-stage}
The denominators \(\Gamma(A_B),\Gamma(A_R),\Gamma(A_D)\) are finite and strictly positive, and the first-stage scores are nondegenerate on the images of \(\Pi^{1/2}A_B\), \(\Pi^{1/2}A_R\), and \(\Pi^{1/2}A_D\).
\end{assumption}

\begin{assumption}[Nondegenerate strict-extension direction]\label{ass:strict-direction}
For strict-extension claims there exists \(v\in\mathcal D\cap\mathcal R^\perp\) such that \(\Pi^{1/2}v\ne0\) and \(\langle \Pi^{1/2}A_Rv,q_{A_R}\rangle\ne0\).
\end{assumption}

\section{Main results}
\begin{theorem}[Theorem 1: filtered population moment decomposition]\label{thm:filtered-decomposition}
Under Assumptions~\ref{ass:finite-linear}--\ref{ass:first-stage}, for each \(A\in\{A_B,A_R,A_D\}\),
\[
  \beta(A)=\beta_{\exp}
  \quad\Longleftrightarrow\quad
  \langle L_A(\epsilon),q_A\rangle=0,
\]
where \(L_{A_B}(\epsilon)=\Pi^{1/2}\epsilon\), \(L_{A_R}(\epsilon)=\Pi^{1/2}A_R\epsilon\), and \(L_{A_D}(\epsilon)=\Pi^{1/2}A_D\epsilon\).  If the first-stage score \(q_A\) is nondegenerate on \(\operatorname{Range}(\Pi^{1/2}A)\), then \(L_A(\epsilon)=0\) is equivalent to consistency of the \(A\)-filtered population IV moment for \(\beta_{\exp}\).
\end{theorem}
\paragraph{Why this is new and worth studying.}
(a) This is routine algebra conditional on Assumption~\ref{ass:moment-expansion}, but it aligns the notation that \citet{BorusyakHullJaravel2022REStud}, \citet{BorusyakHull2023Econometrica}, and \citet{BorusyakHullJaravel2025ECTJ} state separately.  The epsilon-gap is that those papers do not put canonical, recentered, and design-filtered population moments into a common \(\Pi^{1/2}A\epsilon\) diagnostic.  (b) The concrete consumer is the formula-IV diagnostic checklist in \citet{BorusyakHullJaravel2025ECTJ}: Theorem~\ref{thm:filtered-decomposition} supplies the algebraic object that the proposed flagship conjecture will classify sharply.

\begin{conjecture}[Conjecture 1: leakage-null strict-extension frontier (NOT YET PROVED)]\label{conj:leakage-null-frontier}
Under Assumptions~\ref{ass:finite-linear}--\ref{ass:first-stage}, the sharp consistency regions of canonical Bartik, Borusyak-Hull recentered IV, and design-based formula IV are exactly
\[
  \mathcal N_B=\ker(\Pi^{1/2}),\qquad
  \mathcal N_R=\ker(\Pi^{1/2}A_R),\qquad
  \mathcal N_D=\ker(\Pi^{1/2}A_D).
\]
Moreover \(\mathcal N_B\subseteq\mathcal N_R\subseteq\mathcal N_D\), and, under Assumption~\ref{ass:strict-direction}, the second inclusion is strict on an open set of finite-support formula-IV designs.  For any leakage vector \(\epsilon\in\mathcal N_D\setminus\mathcal N_R\), design-based formula IV identifies \(\beta_{\exp}\) while the recentered-IV population limit differs from \(\beta_{\exp}\) by \(\langle \Pi^{1/2}A_R\epsilon,q_{A_R}\rangle/\Gamma(A_R)\).
\end{conjecture}
\paragraph{Why this is new and worth studying.}
(a) The closest published result is \citet{BorusyakHull2023Econometrica}: recentering subtracts the counterfactual average exposure, but the paper does not state an iff theorem separating leakage directions killed by recentering from those killed only by richer design conditioning.  The epsilon-gap is the strict set \(\mathcal N_D\setminus\mathcal N_R\), a projection-null object built from \(\Pi,A_R,A_D,\epsilon\).  (b) Resolving this gives \citet{BorusyakHullJaravel2025ECTJ}'s design-based formula-IV program a concrete decision rule for when the practitioner must use design conditioning rather than a recentered instrument alone.

\begin{conjecture}[Conjecture 2: nondegenerate finite-support witness (NOT YET PROVED)]\label{conj:witness}
Assume Assumptions~\ref{ass:finite-linear}--\ref{ass:strict-direction}.  There exists a finite-support shock design, an exposure matrix \(S\), and an open neighborhood of leakage vectors around some \(v\in\mathcal D\cap\mathcal R^\perp\) such that \(L_D(v)=0\), \(L_R(v)\ne0\), and \(L_B(v)\ne0\).  In that neighborhood, the design-IV limit remains \(\beta_{\exp}\) to first order along \(\mathcal D\), while the recentered and canonical limits move linearly with slopes \(\Gamma(A_R)^{-1}\langle \Pi^{1/2}A_Rv,q_{A_R}\rangle\) and \(\Gamma(A_B)^{-1}\langle \Pi^{1/2}v,q_{A_B}\rangle\).
\end{conjecture}
\paragraph{Why this is new and worth studying.}
(a) \citet{AdaoKolesarMorales2019QJE} show that exposure geometry changes inference, and \citet{GoldsmithPinkhamSorkinSwift2020AER} show that Bartik weights reveal which shares drive the estimate, but neither gives an open-class witness where recentering and design conditioning identify different structural limits.  The epsilon-gap is nondegeneracy: the witness is not a collapsed \(Y\equiv0\) example, but an open finite-support class with a nonzero first stage and nonzero \(\Pi\)-loaded leakage.  (b) The concrete consumer is the immigration-style caution in \citet{JaegerRuistStuhler2018NBER}: a finite-support witness would show when persistent or structured shock exposure is a design-conditioning problem rather than merely a warning about conventional Bartik interpretation.

\section{Sanity-check corollaries}
\begin{corollary}[Zero-leakage benchmark]\label{cor:zero-leakage}
Under Assumptions~\ref{ass:finite-linear}--\ref{ass:first-stage}, if \(\epsilon=0\), then \(L_B(\epsilon)=L_R(\epsilon)=L_D(\epsilon)=0\) and all three population moments reduce to \(\beta_{\exp}\).  This recovers the no-leakage benchmark behind the shock-level orthogonality of \citet{BorusyakHullJaravel2022REStud} and the correctly specified recentering benchmark of \citet{BorusyakHull2023Econometrica}.
\end{corollary}

\begin{corollary}[No strict extension when the spans coincide]\label{cor:spans-coincide}
Under Assumptions~\ref{ass:finite-linear}--\ref{ass:first-stage}, if \(\mathcal R=\mathcal D\), then \(A_R=A_D\), \(L_R=L_D\), and the recentered and design-IV consistency regions coincide.  Hence Conjecture~\ref{conj:leakage-null-frontier} can only be a strict extension when design conditioning removes at least one leakage direction that recentering does not remove.
\end{corollary}

\begin{corollary}[Exposure aggregation can hide leakage]\label{cor:pi-null}
Under Assumptions~\ref{ass:finite-linear}--\ref{ass:first-stage}, if \(\epsilon\in\ker(\Pi^{1/2})\), then all three filters are consistent even if \(\epsilon\ne0\).  The reduction is \(L_A(\epsilon)=\Pi^{1/2}A\epsilon=0\) whenever the relevant filtered leakage is annihilated by the exposure covariance, matching the idea that only \(\Pi\)-loaded leakage enters the formula-IV moment.
\end{corollary}

\section{Proof-sketch route}
The verification route is finite-dimensional linear IV algebra.  First write each filtered instrument as \(SAG\), residualize observation-level controls by \(M_W\), and aggregate the moment to shock space to obtain a numerator of the form \(\langle \Pi^{1/2}A\epsilon,q_A\rangle\).  Theorem~\ref{thm:filtered-decomposition} then follows by dividing by the positive first-stage denominator.  For Conjecture~\ref{conj:leakage-null-frontier}, the hard step is to prove that the moment expansion is sharp in both directions: every nonzero \(\Pi^{1/2}A\epsilon\) can be paired with a nondegenerate first-stage score that makes the bias visible, while zero residuals are annihilated for all admissible scores.  For Conjecture~\ref{conj:witness}, the boundary case should reduce to a three-shock construction with \(\mathcal R\) one-dimensional, \(\mathcal D\) two-dimensional, and \(\Pi\) full rank on \(\mathcal D\cap\mathcal R^\perp\), followed by an open-neighborhood argument using continuity of the IV numerator and denominator.

\section{Honest scope flags}
Theorem~\ref{thm:filtered-decomposition} is labelled as a theorem only conditional on the explicit moment expansion in Assumption~\ref{ass:moment-expansion}; it is not claiming that the expansion has already been derived from every primitive formula-IV design in the published papers.  Conjectures~\ref{conj:leakage-null-frontier} and~\ref{conj:witness} are the open novelty kernels and are not yet proved.  The main uncertainty is whether \(\epsilon\) should be represented directly in shock space or after an \(\Omega^{1/2}\) whitening; the proposal fixes the latter convention in Section~6 to keep dimensions explicit.  The burned angle was the broader spectral \(\lambda_{\max}(\Omega\Pi)\le\lambda_\star\) triple-equivalence threshold; this pivot withdraws that claim and replaces it with a projection-null strict-extension theorem centered on leakage directions.

\section{Iteration trail}
\begin{itemize}[leftmargin=*]
\item v1 (pivot) --- initial draft for angle 1; prior verdict: none.  Replaced the rejected spectral-threshold angle with a structurally different leakage-null projection frontier for recentered versus design-based formula IV.
\end{itemize}

\section*{Bibliography}
\begin{thebibliography}{99}
\bibitem[Borusyak, Hull, and Jaravel(2022)]{BorusyakHullJaravel2022REStud}
Borusyak, Kirill, Peter Hull, and Xavier Jaravel. 2022.
``Quasi-Experimental Shift-Share Research Designs.''
\emph{Review of Economic Studies} 89(1): 181--213.
doi:10.1093/restud/rdab030.

\bibitem[Borusyak and Hull(2023)]{BorusyakHull2023Econometrica}
Borusyak, Kirill, and Peter Hull. 2023.
``Nonrandom Exposure to Exogenous Shocks.''
\emph{Econometrica} 91(6): 2155--2185.
doi:10.3982/ECTA19367.

\bibitem[Borusyak, Hull, and Jaravel(2025)]{BorusyakHullJaravel2025ECTJ}
Borusyak, Kirill, Peter Hull, and Xavier Jaravel. 2025.
``Design-Based Identification with Formula Instruments: A Review.''
\emph{The Econometrics Journal} 28(1): 83--108.
doi:10.1093/ectj/utae003.

\bibitem[Goldsmith-Pinkham, Sorkin, and Swift(2020)]{GoldsmithPinkhamSorkinSwift2020AER}
Goldsmith-Pinkham, Paul, Isaac Sorkin, and Henry Swift. 2020.
``Bartik Instruments: What, When, Why, and How.''
\emph{American Economic Review} 110(8): 2586--2624.
doi:10.1257/aer.20181047.

\bibitem[Adao, Kolesar, and Morales(2019)]{AdaoKolesarMorales2019QJE}
Adao, Rodrigo, Michal Kolesar, and Eduardo Morales. 2019.
``Shift-Share Designs: Theory and Inference.''
\emph{Quarterly Journal of Economics} 134(4): 1949--2010.
doi:10.1093/qje/qjz025.

\bibitem[Jaeger, Ruist, and Stuhler(2018)]{JaegerRuistStuhler2018NBER}
Jaeger, David A., Joakim Ruist, and Jan Stuhler. 2018.
``Shift-Share Instruments and the Impact of Immigration.''
NBER Working Paper 24285.
doi:10.3386/w24285.
\end{thebibliography}

\end{document}
